\begin{abstract}
Logit-lens-style methods interpret hidden states by projecting them through the
LM head and inspecting vocabulary logits, token rankings, or logit differences.
After a token or contrast is selected, however, the inspected quantity is still
a single dense dot product with an LM-head row or linear combination of rows.
This score does not show which LM-head directions support it, oppose it, or
remain in the reconstruction residual. We introduce Sparse Readout Prism (SRP),
a sparse feature decomposition for logit-lens readouts. SRP factorizes the LM
head into shared \emph{readout features} and sparse coefficients for
output-token rows. This factorization first projects each hidden state onto the
shared readout-feature basis; for any selected token logit, logit difference, or
fixed LM-head row combination, it then decomposes the score into signed feature
contributions in original logit units plus an explicit residual.

This fixed basis lets SRP compute hidden-state feature projections once and
reuse them across different logits or contrasts. SRP also composes with
residual-stream attribution and identifies candidate feature directions for
constrained edits to selected readout scores. Across an eight-model suite, SRP
preserves held-out contrast signs on at least \(91\%\) of logit differences,
and on Qwen3.5-2B outperforms nearest-row ridge, matched-rank PCA, and
shuffled/random sparse-code controls. We interpret feature-level decompositions
only when reconstruction preserves the contrast sign and has low error.
\end{abstract}
